%% =====================================================================
%% PART 3 WRITING SKELETON — empty, with instructions
%% Paste over the Part 3 block in main.tex. Write where each "WRITE:" sits,
%% then delete every instruction comment before submitting.
%%
%% Target: ~400 words total. Professional, constructive tone.
%% Judge the WORK, not the person. Every judgement needs one concrete example.
%%
%% Video reviewed: Surbhi Chauhan, "Understanding E-commerce Customers:
%% Predicting Purchase Intention and Customer Churn". Duration 5:52.
%% =====================================================================

\section{Part 3: Peer Review}
\label{sec:part3}

\paragraph{Video reviewed.} Surbhi Chauhan, ``Understanding E-commerce Customers: Predicting Purchase Intention and Customer Churn'' (5:52).

%% ---------------------------------------------------------------------
\subsection{Summary}
%% TARGET: ~60 words. No judgement yet — just what the video presents.
%% Answer: which two business questions; which datasets and sizes;
%% which two models; what she recommends at the end.
%% (From the slides: 12,330 online shopping sessions; 3,941 customer records;
%%  Logistic Regression vs Random Forest; recommends Random Forest as the
%%  starting model for both tasks. CHECK these against what she says.)

% WRITE: summary here.

%% ---------------------------------------------------------------------
\subsection{Soundness and Significance}
%% TARGET: ~90 words. One strength + one gap, each with a timestamp.
%% Prompts — answer from the video, not from the slides alone:
%%   - Are the claims supported by the results she shows?
%%   - Does she state a baseline, how the models were validated
%%     (single split? cross-validation?), or the class balance?
%%   - Does she distinguish association from causation? (a slide does —
%%     check whether she says it too)
%%   - Is "no single best model for every objective" justified by her numbers
%%     (RF ROC-AUC 0.928/0.952 vs LR catching 79.6% of purchasers,
%%      89.7% of churners)?
%% FORMAT: "At <mm:ss> she ___, which ___." Then the gap, same shape.

% WRITE: soundness and significance here.

%% ---------------------------------------------------------------------
\subsection{Executive-Level Insights}
%% TARGET: ~90 words.
%% The test: after watching, would a non-technical manager know what to DO?
%% Prompts:
%%   - The Identify / Detect / Target / Monitor framework — does it give
%%     actions, or just categories?
%%   - Are ROC-AUC and F1 explained in plain language, or stated as numbers?
%%     (the rubric asks for jargon-free)
%%   - Is there any cost, volume or dollar framing — how many customers
%%     would be contacted, what a miss costs?
%%   - Does the recommendation come with a condition
%%     ("validate on current organisational data") — credit that if so.

% WRITE: executive-level insights here.

%% ---------------------------------------------------------------------
\subsection{Presentation Quality}
%% TARGET: ~70 words. Observable facts first, then delivery.
%% Facts you can verify:
%%   - Length 5:52 against the 3--5 minute requirement.
%%   - Model-comparison slide carries a leftover drafting line
%%     ("Your notebook supports these exact results.") and a stray "." bullet.
%%   - That slide's bar chart appears to show four groups but only three
%%     x-axis labels — check whether the fourth is cut off.
%%   - Final slide reads "FINAL TAKEAWAYVISIT" (missing space/line break).
%% Delivery (judge from the video): pace, clarity, audio, whether she
%% reads the slides verbatim, whether slides are readable at speed.
%% Note the slip as a proofreading issue. Do not speculate about its origin.

% WRITE: presentation quality here.

%% ---------------------------------------------------------------------
\subsection{Suggestions}
%% TARGET: ~80 words. Three or four, each naming a specific slide or moment.
%% Specific beats general: "cut the two dataset-description bullets on the
%% second slide to get under five minutes" >> "make it shorter".
%% Aim for: one on length, one on jargon/plain language, one on the slide
%% errors, one on evidence (baseline / validation / cost framing).
%% End with one sentence naming a genuine strength.

% WRITE: suggestions and closing line here.

%% =====================================================================
%% BEFORE MOVING ON
%%  [ ] ~400 words (paste into a word counter)
%%  [ ] Video named exactly, with duration
%%  [ ] Each of the three criteria has a concrete example, ideally timestamped
%%  [ ] 3--4 actionable suggestions
%%  [ ] Professional tone throughout; positive closing line
%%  [ ] All instruction comments deleted
%% =====================================================================
